% !TEX program = lualatex
% ECON 201, Worksheet 9: Choosing the Price and Quantity That Maximize Profit.
% Compile with LuaLaTeX so the PDF is tagged (see syllabus/Econ201-Syllabus.tex
% for the same setup). Build from the course root:
%   latexmk -lualatex -cd worksheets/worksheet09.tex && latexmk -c -cd worksheets/worksheet09.tex
% Every number here is computed and printed by slides/figures/make_lecture9.py
% so the worksheet and the deck quote the same figures (Boba Break).
% A short sheet for the end of class: marginal cost and marginal revenue at
% one point each, worked as intuition rather than a table.
\DocumentMetadata{
  pdfstandard=UA-2,
  pdfversion=2.0,
  lang=en-US,
  tagging=on
}
\documentclass[11pt]{article}

\usepackage[letterpaper, margin=0.75in, top=0.5in, bottom=0.45in]{geometry}
\usepackage{fontspec}
\setmainfont{Fira Sans}
\usepackage{amsmath}
\usepackage{amssymb}
\usepackage{graphicx}
\usepackage{booktabs}
\usepackage{array}
\usepackage{xcolor}
\usepackage{needspace}
\usepackage[most]{tcolorbox}
\definecolor{cream}{HTML}{FDF3EA}
\definecolor{accent}{HTML}{BF5700}
\usepackage{titlesec}
\titleformat{\section}{\large\bfseries\color{accent}}{}{0pt}{}
\titlespacing*{\section}{0pt}{4pt}{2pt}
\usepackage{hyperref}
\hypersetup{pdftitle={ECON 201 Worksheet: Choosing the Price and Quantity That Maximize Profit}, pdfauthor={Div Bhagia}, hidelinks}
\pagestyle{empty}
\setlength{\parindent}{0pt}
\setlength{\parskip}{2pt}

% Ruled answer space: n lines. Plain TeX loop; pgffor's \foreach breaks the
% enumerate counter when used inside an item.
\newcount\answerlinecount
\newcommand{\answerlines}[1]{%
  \par\vspace{1pt}\answerlinecount=#1
  \loop\ifnum\answerlinecount>0
    \rule{\linewidth}{0.3pt}\par\vspace{5pt}%
    \advance\answerlinecount by -1
  \repeat
  \vspace{-5pt}}

% Ruled grids rather than booktabs tables: students write in the cells, so
% every cell gets visible borders and the blank rows get extra height.
\newcolumntype{C}{>{\centering\arraybackslash}p{2.3cm}}
\newcommand{\blankrow}{\rule{0pt}{1.8em}}

\begin{document}

{\LARGE\bfseries\color{accent} Worksheet: Choosing Price and Quantity}\hfill{\large ECON 201}

% Terms sit on the cream panel, where the brand orange falls below 4.5:1;
% the darker orange keeps WCAG AA contrast (see CLAUDE.md).
\definecolor{accentdark}{HTML}{8F4100}
\newcommand{\term}[1]{\textcolor{accentdark}{\textbf{#1}}}

% Definitions panel list: flush left, compact. Built on the generic list
% environment so the bullet sits at the panel's left edge; enumitem is avoided
% for the tagging reason noted in worksheet06.
\newenvironment{deflist}
  {\begin{list}{\textbullet}{%
     \setlength{\leftmargin}{1.1em}\setlength{\labelwidth}{0.7em}%
     \setlength{\labelsep}{0.4em}\setlength{\itemindent}{0pt}%
     \setlength{\listparindent}{0pt}\setlength{\topsep}{2pt}%
     \setlength{\partopsep}{0pt}\setlength{\parsep}{0pt}%
     \setlength{\itemsep}{2pt}}}
  {\end{list}}

\begin{tcolorbox}[colback=cream, colframe=accent, boxrule=0.6pt, arc=2pt,
  left=6pt, right=6pt, top=4pt, bottom=4pt, title=Definitions,
  fonttitle=\bfseries, coltitle=white]
\begin{deflist}
\item \term{Profit}: $\text{Profit} = R - C(Q)$, where revenue is $R = P \times Q$.
\item \term{Marginal cost}: the increase in total cost when one additional unit
  is produced.
\item \term{Marginal revenue}: the increase in revenue when one additional unit
  is sold. To sell one more, the firm lowers its price on every unit, so
  \begin{deflist}
  \item $\text{MR} = $ price of the extra unit $-$ the loss on the units it
    was already selling.
  \end{deflist}
\item \term{Profit-maximizing quantity}: where $\text{MR} = \text{MC}$.
  \begin{deflist}
  \item While $\text{MR} > \text{MC}$, one more unit adds to profit. Once
    $\text{MR} < \text{MC}$, one more unit lowers it.
  \end{deflist}
\end{deflist}
\end{tcolorbox}

\needspace{6\baselineskip}\section{Boba Break}

Boba Break is a stall at the Tuesday farmers market on campus. It pays \$100 a
day for the stall and \$4 to make each cup, so $C(Q) = 100 + 4Q$. Its demand
curve is $P = 12 - Q/10$: the more cups Boba Break wants to sell, the lower the
price it has to charge.

\needspace{5\baselineskip}\section{1. Marginal cost}

How much does one more cup add to Boba Break's cost? Does it matter how many
cups it is already selling?
\answerlines{2}

\needspace{9\baselineskip}\section{2. Marginal revenue at 30 cups}

(a) What price does Boba Break have to charge to sell 30 cups? What about 31
cups?
\answerlines{1}

(b) What is its revenue at 30 cups, and at 31 cups?
\answerlines{1}

(c) What is the marginal revenue? Should Boba Break sell the 31st cup?
\answerlines{1}

\needspace{7\baselineskip}\section{3. Marginal revenue at 50 cups}

Now Boba Break sells 50 cups at \$7, and a 51st cup would sell for \$6.90. What
is the marginal revenue? Should it sell the 51st cup?
\answerlines{2}

\needspace{6\baselineskip}\section{4. The best quantity}

Somewhere between 30 and 50 cups, marginal revenue falls below \$4. Where
should Boba Break stop, and what price does it charge?
\answerlines{2}

\end{document}
